\section{n2. Full conductor-envelope equivalence when the entire normalized boundary is bijective; explicit radicial and primary-conductor criteria.}
\textit{Supports: G331, n1.}

New conductor theorem, relative to ENTIRE qualified G331 and own n1. Let R be excellent reduced local with finite normalization S and conductor c. Set X=V\_a(R) and Xtilde=V\_(aS)(S), using every closed normalization centre and EVERY semivaluation support. Suppose restriction on the conductor quotient boundaries is injective. More precisely, the boundary of (S/c)\_red maps injectively to that of (R/c)\_red, with the quotient of the extended gauge, omitting empty zero-dimensional sections. Then rho:Xtilde-\textgreater{}X is a homeomorphism. Indeed G331 grants surjectivity and compact Hausdorff quotient topology. At finite conductor value every fibre is singleton. At infinite value its ENTIRE fibre is precisely the conductor-quotient restriction fibre. The assumed injectivity handles all those fibres. Compact bijectivity gives the homeomorphism. This is a sufficient condition, not an asserted necessity for radiciality of an algebraic conductor map.

Here are two algebraically checkable sufficient cases. First, if R/c has dimension zero, c contains a power of m, so every centred semivaluation has finite conductor value. No quotient boundary remains; multiple normalization centres are allowed. Second, let B=(R/c)\_red and T=(S/c)\_red. Assume each prime p of B has exactly one prime P of T above it, and the finite field extension Frac(T/P)/Frac(B/p) is purely inseparable (in characteristic zero this means equal). These are the usual finite radicial conditions, stated explicitly. The reduction map B-\textgreater{}T is injective: nilpotence upstairs of an element of R/c is already nilpotence in its injected subring. For a quotient semivaluation with support p, any lift has support P, uniquely. After quotienting the supports, its value on any nonzero field element x is forced by w(x)=v(x\textasciicircum{}(p\textasciicircum{}e))/p\textasciicircum{}e when an appropriate power lies in the old field; in characteristic zero it is already the same field. This proves uniqueness on every support, not just generic valuations. Existence is G331 finite lifting. Thus the quotient boundary is injective. No global common Frobenius exponent, perfection, F-finiteness or characteristic-zero restriction is required for this uniqueness proof. We do not invent a real field valuation through a nonzero support.

Because EVERY rho-fibre is singleton, G331's FULL positive C and D descent implies order bijections rho*:C\_R-\textgreater{}C\_S and rho*:D\_R-\textgreater{}D\_S. On C this is an isometry, since rho is surjective; continuity in either direction follows from the homeomorphism. Thus, for ANY obstacle for which the raw supremum is interpreted pointwise, the completed order bijection gives p\_S(rho*f)=rho*p\_R(f). For finite continuous signed obstacles the raw suprema are finite by n1. The continuous-envelope property is EQUIVALENT on R and S: every continuous source obstacle is the pullback of a continuous target obstacle, and an envelope belongs to C on one side iff it does on the other. The descended continuous envelope is independently geometric by positive descent, not by an unsupported assertion that a pointwise supremum is positive.

The canonical degree-m ideals satisfy A\_m\textasciicircum{}R(f)=A\_m\textasciicircum{}S(rho*f) intersect R, directly from surjective restriction and ALL defining inequalities. Equality after extension to S is NOT asserted. When CE holds, n1 gives uniform convergence in EVERY downstairs canonical degree. Each such profile has an ACTUAL negative effective Cartier presentation -E\_(A\_m)/m on its ACTUAL primary blowup; its degree on EVERY contracted integral curve is positive by ampleness. Consequently the uniform geometric envelope has actual positive approximants on the required model system. A continuous envelope need not itself have a single-model determination. Neither a source O(1/m) bound nor source finite generation is automatically inherited through contraction; n3 establishes an explicit rate and failure of finite generation for a particular family by a different argument.

The exact primary-gauge cocycle commutes with rho since w(aS)=v(a). For a new gauge b, d(v)=v(b), source and target divide by the same bounded positive denominator. Canonical ideals remain unchanged when the obstacle is transported homogeneously. Completion follows G331 and n1. For decreasing continuous obstacles with continuous envelopes, the actual D rule of n1 transports as well, with partial bottom and bottom retained. This theorem supplies envelope descent only under the stated singleton-fibre criterion; it does not settle arbitrary nonnormal gluing along a positive-dimensional conductor quotient.

Consequences: any excellent reduced local germ whose conductor is primary has CE exactly when all of its normal local centres have CE. In particular, the dimension-two normal theorem of n1 gives CE for every such germ whose normal centres all have dimension at most two. Under the finite radicial conductor-quotient condition the same equivalence holds even when the infinite-conductor boundary is NONEMPTY. If S is a normal local surface it has CE in any equal/mixed characteristic by n1; if its centres are smooth affine closed germs over an algebraically closed characteristic-zero field it has CE in every dimension by n1. These are narrowed conclusions with explicit source hypotheses, not added premises for unrestricted ORIGINAL P12.3.

The point where the unrestricted descent problem remains is now precise: for a quotient boundary with nonsingleton fibres there is no automatic constancy of a source greatest minorant. This note neither takes a fibrewise minimum as psh nor asserts arbitrary finite-map envelope compatibility. Products, energy and full finite operators remain separate.

This is a mathematical proof using G331 only for its qualified boundary/cone/descent results and n1 for envelopes. No new published premise is needed for the radicial uniqueness argument. Artifact n2.md, expanded n2.tex and the standalone retain the exact theorem; rerun python assemble.py; python verify.py. Finite hash/serialization success does not verify the theorem.

\subsection*{Expanded conductor theorem}
With $\mathfrak c=(R:S)$, assume the reduced conductor quotient boundary
restriction is injective. Then
\[
\rho:V_{\mathfrak aS}(S)\xrightarrow{\sim}V_{\mathfrak a}(R),\quad
\rho^*C_R=C_S,\quad\rho^*D_R=D_S,
\]
\[
\rho^*p_R(f)=p_S(\rho^*f),\qquad
\operatorname{CE}(R)\Longleftrightarrow\operatorname{CE}(S),\qquad
A_m^R=A_m^S\cap R.
\]
This condition includes the entire infinite-conductor support stratum.
The last identity is contraction only, not an extension identity.

